\documentclass[10pt]{article}
\usepackage[margin=20mm]{geometry}
\usepackage{amsmath,amssymb,fontspec,unicode-math,xurl,hyperref}
\IfFontExistsTF{DejaVu Serif}{\setmainfont{DejaVu Serif}}{\setmainfont{Latin Modern Roman}}
\setmathfont{latinmodern-math.otf}
\setlength{\parindent}{0pt}\setlength{\parskip}{6pt}\sloppy
\begin{document}
Abstract: The complete published paper, its appendix and the supplied principal model, training and evaluation files were inspected. The claimed mathematical derivation assumes a Lipschitz gradient for an \(\ell_1\) norm and contains an incorrect normal-equation coefficient. The current release omits reconstruction loss from optimization, does not feed reconstructed backgrounds into successive segmentation updates, and creates state-space blocks afresh inside forward calls. These are substantive failures of the asserted strict correspondence between the objective, network and experimental implementation. The empirical benchmark results are not thereby proved false; they require an identified historical program and corrected mechanism evidence. Scientific adjudication remains pending.\par
Author: AI-assisted technical review\par
Date:\par
Title: Technical Verification of RUN: Reversible Unfolding Network for Concealed Object Segmentation\par
\# Identity and scope\par
This report reviews Chunming He, Rihan Zhang, Fengyang Xiao, Chengyu Fang, Longxiang Tang, Yulun Zhang, Linghe Kong, Deng-Ping Fan, Kai Li and Sina Farsiu, ICML 2025, PMLR 267, I25-43768 (0399), OpenReview `tiYNbKwL7T`. All 12 physical pages were read, including references, Algorithm S1 and Appendix A. All Figures 1–6 were actually viewed. Printed and physical page numbers agree. The acquired repository is `ChunmingHe/RUN`, commit `878208d5154577075fe15``156866970f2a57c2fb8`; all acquired file hashes were checked.\par
Complete static reading covered `Network.py`, `get\_m\_hat.py`, `Train.py`, `Test.py`, `data\_val.py`, `utils.py`, `eval\_functions.py` and the README. In `Modules.py`, the state-space wrapper, selective Mamba wrapper, REU blocks and complete REM11 were inspected (lines 1–68, 297–338, 920–1045). The remaining low-level state-space internals and unrelated utility/backbone implementations were not fully audited. No author code was executed or imported. Independent scalar calculations and static AST checks, including their command outputs, are preserved. The source is an acquired subset of a public release; neither its correspondence to historical experiments nor inaccessible separate OpenReview attachments was established.\par
\# Supported findings\par
Problem 1 (The asserted optimization derivation is invalid for the specified objective).  Equation (4), p. 3, explicitly defines \(S\) as an \(\ell_1\) norm. Equations (10)–(12), p. 4, replace it by a quadratic Taylor model whose justification assumes that \(\nabla S\) is Lipschitz continuous. For the scalar restriction \(S(r)=|r|\), the derivative is undefined at zero, and for \(r=\varepsilon\) and \(s=-\varepsilon\), the derivative difference divided by \(|r-s|\) equals \(1/\varepsilon\). No finite global gradient-Lipschitz constant exists.\par
This obstruction is directly relevant to the proposed initialization: zero masks and the piecewise mapping in Equation (5) permit zero residuals. Choosing zero as a subgradient at zero does not supply a differentiable function or a valid tangent quadratic upper model: \(|r|>Lr^2/2\) for sufficiently small nonzero \(r\) and any finite \(L\). An additive constant that makes a quadratic a loose bound cannot restore the missing first-order equality required by the stated Taylor argument. A deliberately smoothed penalty could fix the premise, but its definition, smoothing scale, approximation error and resulting update are absent.\par
Even if Equation (12) is taken simply as a chosen quadratic surrogate, its printed solution is wrong. Interpret image products charitably as pointwise operations and consider one scalar pixel, freezing \(w_k,\widetilde m_k,R_{k-1}\) as stipulated. The coefficient of \(\widehat m\) in the normal equation must be \[C^2+\mu+\alpha L_S w_k^2,\] whereas Equation (13) sets \(Q_a=C^2+\mu+L_Sw_k^2\), omitting \(\alpha\). The numerator retains \(\alpha\), so a uniform redefinition of \(L_S\) does not remove the discrepancy. For \(C=1\), \(B=0\), \(\mu=1\), \(\alpha=2\), \(L_S=10\), \(m_{k-1}=0.2\), \(\widetilde m_k=0.1\), \(w_k=w_{k-1}=1\), and \(R_{k-1}=0.1\) with gradient one, Equation (12) has \[J(m)=\tfrac12(1-m)^2+\tfrac12(m-0.2)^2+10(m-0.1)^2.\] Its unique minimizer is \(3.2/22\approx0.145455\). Equation (13) instead gives \(3.2/12\approx0.266667\), where \(J'(m)=2.666667\ne0\). This scalar case already disproves the displayed closed-form claim. If one scalar mask is shared across RGB channels, the data-term coefficient additionally needs the appropriate channel sum; averaging independent channelwise solutions is not generally the same minimization.\par
The supplied `get\_m\_hat.py`, lines 39–49, differs further: its denominator adds a constant \(L_S\) without \(w_k^2\) or \(\alpha\), while its numerator retains their product. The constructor’s \(L_s\) argument is stored but the formula resets \(L_s\) to an all-ones tensor. Those choices may define a useful learned architecture, but they do not verify the claimed exact optimization correspondence.\par
This is Major because the paper emphasizes mathematically strict connections in Figure 3 and explicit theoretical guarantees on p. 8. It is a derivation failure, not a proof that the learned model must perform poorly. Use a valid nonsmooth proximal update or state and analyze a smooth surrogate; rederive the normal equations; define all pointwise operations, channel reductions and initial states, including \(M_{-1}\) needed at \(k=1\). If arbitrary learned blocks replace proximal maps and gradients, identify which guarantees survive and under what assumptions. Otherwise describe the method as optimization-inspired and remove unsupported convergence or optimality wording. The background scalar update in Equation (17) was checked and does solve its stated quadratic surrogate; that local correctness does not repair the mask derivation.\par
Problem 2 (The supplied training and inference paths omit the claimed joint unfolding mechanism).  Equation (22), p. 5, includes stage-wise reconstruction loss, and Figures 2–3 and Algorithm S1 require a refined background to influence the next segmentation stage. In the advertised `Train.py`, lines 80–82, a reconstruction quantity is computed, but line 83 sets the differentiated objective to `loss\_init + loss\_final + loss\_edge + 2 * ual\_loss`. The reconstruction quantity is absent; it is only logged afterward.\par
The extra uncertainty loss is defined in `utils.py`, lines 29–33, as the mean of \(1-(2\sigma(z)-1)^2\), rather than the reconstruction term in Equation (22). Moreover, in the default `Network.forward`, `Network.py` lines 646–682, the mask pipeline and selective Mamba branch modify only the coarse returned mask after `REM11` has already computed all four masks. Reconstructed backgrounds are then computed at lines 669–680, with no call returning them to another segmentation update. `REM11`, lines 998–1025 of `Modules.py`, was read to confirm the ordering and independent returned tensors. `Test.py` line 55 selects return element four, the fine mask `f1` computed before those modifications. Thus the reconstructed background branch is absent from the final forward dependency and, with the reconstruction term omitted, its own parameters have no path to the optimized loss. Auxiliary supervision on the modified coarse mask can still affect shared encoder training; this review does not claim that every added branch has no training effect.\par
A further implementation defect affects the central RSS component. `Network.py` line 655 creates a local `Selectivemamba()` on each call; its own forward method creates local `VSSBlock` instances at `Modules.py` lines 324 and 329. These are newly initialized module objects, not registered persistent children of the model used by the optimizer at `Train.py` line 214 or by checkpoint loading. They are therefore not learned and restored as the paper’s fixed trained RSS blocks. Calling `model.eval()` cannot turn newly created, unregistered objects into persistent evaluated submodules. This statement follows from object construction and registration, not an executed training trace.\par
These discrepancies constitute a Major obstacle to reproducing the central joint reconstruction/segmentation mechanism. A different historical program or accidentally released development branch could explain them; the published tables are not proved fabricated or numerically wrong. Publish the source and configurations that generated them, register state-space modules in constructors, make every intended background-to-mask dependency explicit, and include the intended reconstruction term in the training objective. Recompute the mechanism ablations with matched checkpoints and verify gradients reach the reconstruction branch. The present loss and graph cannot substantiate the claimed reconstruction-guided causal explanation merely because an intermediate image is displayed in Figure 2.\par
Problem 3 (The claimed experimental breadth and some universal comparisons exceed the available results).  Contribution (3), p. 3, claims validation on five COS tasks as well as salient object detection. The complete published PDF supplies comparative results for COD, polyp segmentation, tubular segmentation and transparent objects. Appendix A gives a CDS2K defect-detection protocol but no corresponding result table or figure, and no salient-object evaluation is provided. Those additional validation claims cannot be checked from the acquired paper.\par
Separately, the broad superiority wording on p. 6 needs metric-specific qualification. ResNet50 RUN on CAMO has \(E_\phi=0.868\), below FocusDiff’s \(0.883\), and \(S_\alpha=0.806\), below FSEL’s \(0.816\). RUN’s CORN AUC is \(0.962\), below DSCNet’s \(0.964\). These exceptions coexist with strong gains elsewhere and are not a reversal of every ranking. This Minor reporting finding requires the missing experiments or narrower scope, and wording that states the observed tradeoffs. Separate inaccessible attachments may contain additional results, but their existence and contents were not verified.\par
\# Independent checks and remaining uncertainty\par
The four highlighted numerical gains in Figure 1 were checked against the relevant displayed table values: \(0.747-0.730=0.017\) for COD10K \(F_\beta\), \(0.878-0.871=0.007\) for ETIS structure score, \(0.870-0.859=0.011\) for CORN sensitivity, and \(0.952-0.944=0.008\) for GDD maximum F-score against the strongest tabulated competitor. These calculations verify differences, not experimental provenance. Table 7’s stage progression is compatible with diminishing gains in \(F_\beta\); the four-stage setting is a declared efficiency choice, although wall time, parameter counts and compute-matched stage comparisons are not reported. Small increments such as \(0.001\) in ETIS Dice need variability information before being characterized as robust improvements.\par
Reconstruction alone does not identify a correct mask: for any mask \(m\), choosing \(B=C(1-m)\) makes \(C-Cm-B=0\). The excluded interval in Equation (5) also gives no residual-sparsity penalty to maximally ambiguous masks. Learned semantic priors and supervised mask/edge losses may resolve that ambiguity, but the reconstruction image by itself is not evidence of segmentation correctness. The visual examples demonstrate selected successes; no random sample protocol or pixel-level recalculation was available. The paper’s term “reversible” describes foreground/background reasoning and was not misread as a formal bijection guarantee.\par
The released loader normalizes RGB data and uses unbounded logits on parts of the mask path; these details should be reconciled with the stated \([0,1]\) mask-domain interpretation. The published zero initialization also differs from the learned prior initialization in the release. Current default validation points to CDS2K and selects the best checkpoint by validation MAE; the test script names a path containing a COD10K validation label. These filenames do not prove which data were used historically or establish leakage. Exact split manifests and checkpoint-selection rules are needed. Likewise, the CORN “first 30\% / last 70\%” split requires a specified ordering and subject grouping for replication.\par
The full experimental metric engine is not established by the supplied utilities: the README delegates evaluation to an external MATLAB toolbox, while local utilities include both weighted and thresholded F-score routines. Appendix A’s precise F-score convention, thresholds and averaging should be tied to the actual evaluator. Baseline retraining parity, hidden tuning, raw masks, error bars and the influence of the extra pretrained dehazer in RUN+ were not independently reproduced. The latter is an additional component and does not isolate the benefit of unfolding alone.\par
\# Repair status\par
The nonsmooth-gradient obstruction, missing normal-equation coefficient and scalar background solution were independently checked. Static loss/dependency evidence was verified without author execution. Two Major findings and one Minor finding are supported. Fresh independent acceptance rechecked both Major Problems and the Minor scope/comparison Problem. Scientific adjudication remains pending; human reviews remain empty.\par
\end{document}
